\section{引言}

团簇化量子自旋链为研究局域多重态、量子临界性、隐藏弦关联和边缘自由度之间的
关系提供了可控平台。已有四聚体链研究建立了“孤立单胞多重态—弦关联—纠缠结构—
动力学结构因子”的证据组织方式，并展示了不同耦合区域的低能激发特征
\cite{li2026tetramer}。然而，五聚体单胞包含奇数个自旋 $\frac{1}{2}$，
其低能自由度和 Lieb--Schultz--Mattis 类型约束与偶数自旋单胞不同，因此不能
直接移植四聚体链的相名、弦算符端点或准粒子命名。

本文考虑具有两个胞内耦合 $J_1,J_2$ 和胞间耦合 $J_3$ 的五聚体
Heisenberg 链。我们首先分析孤立五聚体的低能多重态，然后在固定
$\beta=J_3/J_1=0.5$ 的截面上计算基态能量、普通自旋关联、弦关联、
边缘磁化、纠缠谱和动力学结构因子。分析的核心问题包括：

\begin{enumerate}
  \item 负 $\alpha$ 区域是否为普通有隙 Haldane-like 区域；
  \item 正 $\alpha$ 区域内部的两个响应峰是否定义额外相；
  \item $\alpha=0$ 两侧的低能多重态、弦关联和动量结构如何变化；
  \item 负侧的边缘和弦特征能否与临界体关联共存。
\end{enumerate}

当前结果表明，$\alpha=0$ 两侧均具有代数体关联，但负侧同时保留明显更慢的
弦关联衰减和局域边缘响应。我们因此将其保守地称为具有弦与边缘特征的临界区域，
而不提前命名为普通有隙 SPT 相或确定的 gapless SPT 相。

\todo{补充更广泛的五聚体链、无隙 SPT 和团簇自旋链文献。}
